\chapter{Implementation}
\label{chap:Implementation}
\lstset{language=Python}
\captionsetup{width=.8\linewidth}

% REVISION NOTE (target +5 to 6 pages): most of the honest extra length belongs in
% this chapter. Reduce repeated WHY and replace it with HOW: Python-side changes,
% bridge lifecycle, state transitions, concrete translation examples, and shorter,
% better-focused listings.

This chapter describes how the design of the previous chapter was realised in the repository. It explains the staged transition from the original Python refiner to the Haskell version, the supporting bridge and diff-based workflow used to keep the port testable, and the point at which the implementation stopped being a direct translation and became a more Haskell-native rewrite.

\section{Overview of the Solution}

% SUPERVISOR NOTE (p17): this opening is starting to repeat the WHAT and WHY.
% Trim repetition here and spend the recovered space lower down on mechanics,
% examples, and concrete code-path explanation.

The implementation followed a gradual transition from the original Python refiner to a Haskell refiner while keeping the rest of the pipeline recognisable. The main files involved were \texttt{refine\_command.py}, \texttt{spin2test.py}, \texttt{combined\_haskell\_bridge.py}, and the new \texttt{refine\_command.hs}. The objective was not to replace everything at once, but to replace the refinement stage in a way that could be checked end to end against the old behaviour.

A full rewrite from the start would have been risky. The refiner is small enough to look manageable, but it contains many state-dependent branches and many points where template expansion affects the final test file. Reaching the end of a full rewrite before comparing outputs would therefore have made it difficult to localise errors. The staged approach reduced that risk by keeping differential comparison in the loop throughout the port.

The testing and debugging effort was directed toward functional equivalence with the Python generator. The Python output already matched the expected RTEMS test structure, so the main implementation question was whether the Haskell version could preserve that observable behaviour while using a clearer internal structure.

\section{The Implementation}

\subsection{The Bridge}

% SUPERVISOR NOTE (p18): add 1.5 to 2 pages here on HOW the bridge works in practice.
% Show the Python entry point, the message format, process startup/shutdown,
% error handling, and one short logged exchange from INIT through testBody.
% Also explain what changed on the Python side to preserve the historical command
% interface while delegating the core refinement work to Haskell.

During the transition, a bridge layer preserved the historical Python-facing \texttt{command} interface while moving the refinement work behind a subprocess boundary. Python remained responsible for orchestrating the wider pipeline, but Haskell could now be exercised inside that pipeline without forcing an immediate rewrite of every caller.

Communication was carried entirely through standard input and output. The bridge sends configuration data to the Haskell process, forwards annotation lines one at a time, and finally requests the rendered test body. This arrangement is deliberately plain, but it has two practical benefits: the interface is explicit, and the same messages can be logged or replayed when something goes wrong.

\subsection{The Diff Script}

% SUPERVISOR NOTE (p19): keep listings short and stop them breaking across pages.
% If necessary, replace full-file listings with excerpted fragments that only show
% the mismatch being discussed, and move secondary detail into prose.

A constant part of the porting process was checking the generated test files against the ones produced by the original system. At first this was done manually with an online text comparison tool, but that approach became unworkable once a single change could affect many generated files. A small local diff script therefore became part of the working process.

\includecode{Sample of command line output}{Example output from the diff workflow showing a generated file diverging from the reference output. In practice, even a single extra line was worth investigating because small formatting mismatches could indicate a deeper error in template expansion.}{lst:snippet}{commandline.py}

The diff-based approach was intentionally strict. It reported trivial whitespace differences as well as genuinely important mismatches. That occasionally produced noise, but it was still useful because it made it harder for a small divergence to go unnoticed and later turn into a larger behavioural discrepancy. In practice, most debugging started from a diff, then moved backwards through the relevant refinement dictionary entry and the code path that emitted it.

Unit tests were also added around the Haskell refiner, but they were less decisive than the end-to-end diff checks. Unit tests were useful for protocol handling, configuration loading, and small local behaviours. They were less effective for the interactions that actually mattered most, where several pieces of state and several refinement templates combined to produce one incorrect line of generated code.

\subsection{Refine Command}

% SUPERVISOR NOTE (p21): add 2.5 to 3 pages here with a running example. Start from
% one concrete annotation line or short sequence, show the relevant Python template or
% dictionary entry, show the Haskell state transition(s), and end with the emitted C.
% This is the cleanest way to add useful length without padding.

This was the main file to be converted. Its job is to take annotated SPIN output and refine it into concrete test code by matching annotation lines against entries in a model-specific refinement dictionary. Those entries are small templates of C code with placeholders for names, values, and identifiers taken from the trace.

At the beginning of the port, the Python code was translated into Haskell as closely as possible. The intention was to keep behaviour stable by keeping the structure stable. What became clear quite quickly, however, was that this strategy made the Haskell version look like Python written in another syntax. It preserved the awkward parts of the original without gaining much clarity from the new language.

The turning point came in state handling. The program has to remember whether the current annotation is inside a structure, inside a sequence, or part of ordinary top-level output. The Python version manages this with mutable object fields that are updated as each line is processed. That behaviour can be copied in Haskell, but copying it directly would have reproduced the same difficulty of following how state changes over time.

The more Haskell-native approach was to make the state explicit. Instead of mutating fields on one long-lived object, each stage receives a runtime state, transforms it, and returns the next one. At that point the port stopped being a line-by-line translation and became a Haskell implementation that happened to preserve the same outputs. The aim from then on was not structural similarity to the Python source, but behavioural equivalence of the generated tests.

This change is visible in the following snippets.

\paragraph{Python state stored and mutated on the object}
\begin{lstlisting}[language=Python, caption={Python state stored and mutated on the object}, label={lst:python-state-mutation}]
self.inSTRUCT = False
self.inSEQ = False
self.seqForSEQ = ""
self.inName = ""

case [pid, "SEQ", name]:
    self.inSEQ = True
    self.seqForSEQ = ""
    self.inName = name

case [pid, "SCALAR", "_", value]:
    self.seqForSEQ = self.seqForSEQ + " " + value
\end{lstlisting}

In the Python version the state lives on the command object. As each SPIN line is read these values are updated in place, and later branches depend on whatever earlier branches left behind.

\paragraph{Haskell state made explicit}
\begin{lstlisting}[language=Haskell, caption={Explicit runtime state in the Haskell refiner}, label={lst:haskell-runtime-state}]
data RuntimeState = RuntimeState
  { rtInSTRUCT :: Bool
  , rtInSEQ :: Bool
  , rtSeqForSEQ :: Text
  , rtInName :: Text
  , rtDefCode :: [Text]
  , rtDeclCode :: [Text]
  , rtTestCodes :: [(Int, [Text])]
  , rtCurrentId :: Int
  , rtSwitchNo :: Int
  }
\end{lstlisting}

The same information is still present, but it is now grouped into one explicit record. Rather than hidden mutation on an object, each function takes a state and returns the next one.

\paragraph{One SPIN line producing the next state}
% SUPERVISOR NOTE (p21): "Fix these!" Review the line wrapping and caption flow around
% Listings 4.4 and 4.5. If they still crowd the page, split them or shorten them.
\begin{lstlisting}[language=Haskell, caption={Processing one SPIN line by returning a new state}, label={lst:haskell-state-pipeline}]
refineSPINLineMain :: [Text] -> RuntimeState -> RuntimeState
refineSPINLineMain ln rt0 =
  let rt1 = collectPIdsTokens ln rt0
      rt2 = if outputSWITCH (rtFlags rt1) then switchIfRequired ln rt1 else rt1
      rt3 = if annoteComments (rtFlags rt2) then logSPINLine ln rt2 else rt2
  in refineSPINLine ln rt3
\end{lstlisting}

This is the clearest expression of the change in approach. Instead of one long-lived object being updated in place, each stage produces a fresh version of the runtime state before the next stage runs.

\paragraph{Sequence handling as an explicit transition}
\begin{lstlisting}[language=Haskell, caption={Sequence state handled as an explicit transition}, label={lst:haskell-sequence-state}]
_pidTxt:"SEQ":name:[] ->
  rt { rtInSEQ = True, rtSeqForSEQ = "", rtInName = name }

_pidTxt:"SCALAR":"_":value:[] ->
  rt { rtSeqForSEQ = rtSeqForSEQ rt <> " " <> value }
\end{lstlisting}

Sequence handling was where this change felt most useful. In this context a sequence is an ordered run of values collected between \texttt{SEQ} and \texttt{END} and later emitted through a matching \texttt{\_SEQ} template. In the Haskell version, entering that mode and accumulating those values are both visible directly in the returned state.

% LENGTH TARGET (+1 to 1.5 pages): after the state snippets, add one end-to-end
% translation walkthrough with intermediate runtime-state snapshots. That will satisfy
% the request for more Haskell detail better than adding more isolated code fragments.

\section{Summary}

The implementation was successful because it kept verification in the loop while the internal structure changed. The bridge kept the surrounding Python pipeline stable, the diff script exposed regressions quickly, and the Haskell refiner could then be reorganised into an explicit-state design without losing confidence in the generated output.

